\acresetall

\chapter[Trajectory-Aware Knowledge Estimation for Text Dataset Distillation]{Trajectory-Aware Knowledge Estimation \\ for Text Dataset Distillation} \label{ch:50-research03}

\begin{foldable}
  The background on text \ac{DD} (\S\ref{sec:26-background-datasetdistill-textdd}) lays out what a complete method has to do, and \S\ref{sec:27-background-influence} and \S\ref{sec:28-background-ot} pin down the two parts no prior method has solved together: at extreme compression, importance scores read off a single checkpoint are systematically pulled toward noisy examples, and even with corrected scores, independent ranking cannot guarantee that the selected subset covers the weighted distribution.
  This chapter develops \textbf{\acl{TAKE}} (\textbf{\acs{TAKE}}), which scores each sample by integrating its influence over the entire training trajectory rather than at a single point, and then uses those scores as the source measure in a discrete Optimal Transport selection over an \ac{LLM}-generated candidate pool.
  The trajectory-aware score corrects the bias that the single-checkpoint baseline introduces, and the \ac{OT} formulation turns sample-level scoring into a corpus-level selection objective whose output remains human-readable.
\end{foldable}

\section{Introduction} \label{sec:51-research03-introduction}

\begin{foldable} % Text DD
  The systemic costs that motivate dataset compression sharpen at the corpus scale: at extreme compression, the residual budget has to be spent on the samples that contribute the most to learning.
  Yet, adapting image \ac{DD} techniques~\cite{wang2018dataset,zhao2021dataset,zhao2023dataset} to text faces two concrete obstacles: non-differentiable token decoding blocks gradient-based synthesis, and embedding spaces are tightly coupled to specific architectures.
\end{foldable}

\begin{foldable} % Baseline text DD methods --> gap --> desiderata
  Existing text \ac{DD} methods~\cite{sucholutsky2021soft,li2021data,maekawa2023dataset,maekawa2024dilm,tao2024textual} fall short in at least one of the following: at extreme compression, uniform weighting exhausts the budget on uninformative samples; their objectives lack task alignment or global optimization; and their outputs are embeddings, neither auditable nor transferable.
  We survey this landscape in \S\ref{sec:26-background-datasetdistill-textdd} and show that no prior method simultaneously addresses all of the above.
  These gaps motivate three desiderata for TAKE:{
    (1)~\textit{knowledge-based weighting} --- weighting samples by their task-aligned downstream contribution, directing the budget toward the most informative ones;
    (2)~a \textit{dataset-level distillation objective} --- optimizing over the full training distribution rather than batch-by-batch, ensuring global coverage;
    and (3)~\textit{human-readable output} --- a distilled corpus that is directly inspectable, auditable, and transferable across architectures.
  }
\end{foldable}

\begin{table}[ht]
  \centering
  \begin{threeparttable}
    \caption{Landscape of text dataset distillation methods.}
    \label{tab:landscape-textdd-methods}
    \setlength{\tabcolsep}{8pt}
    \begin{tabular}{l c c c c c c}
      \toprule
      \textbf{Gap}         & SLDD         & DDTC         & DDAL         & DiLM         & DaLLME       & \textbf{TAKE} \\
      \midrule
      Task alignment       & $\checkmark$ & $\checkmark$ & $\triangle$  & $\triangle$  & $-$          & $\checkmark$ \\
      Global optimization  & $\triangle$  & $\triangle$  & $\triangle$  & $\triangle$  & $\triangle$  & $\checkmark$ \\
      Interpretability     & $\triangle$  & $-$          & $-$          & $\checkmark$ & $\checkmark$ & $\checkmark$ \\
      Importance weighting & $-$          & $-$          & $-$          & $-$          & $-$          & $\checkmark$ \\
      \bottomrule
      \multicolumn{7}{l}{$-$: absent, $\triangle$: partial, $\checkmark$: addressed}
    \end{tabular}
  \end{threeparttable}
\end{table}

\begin{foldable} % HSB framing and positioning
  The \textit{importance weighting} row in Table~\ref{tab:landscape-textdd-methods} stands out: every prior method leaves it open.
  The natural fix is influence functions~\cite{koh2017understanding}, which quantify each sample's counterfactual effect on the downstream loss.
  Scoring at a single checkpoint, however, introduces the \textit{\acl{HSB}} (\ac{HSB}): at convergence, clean and moderate samples have near-zero gradient by definition, so noisy or hard samples dominate the influence scores, exactly inverting the ranking a robust distilled corpus needs.
  In standard supervised training the bias can often be absorbed by regularization, but in the extreme low-data regime of \ac{DD} it amplifies instead: with only a handful of examples per class to allocate, biased selection actively suppresses the informative samples the surrogate depends on.
  To our knowledge, no prior text \ac{DD} method has identified or corrected for the \ac{HSB}, which is why the table's last row remains uniformly empty.
\end{foldable}

\begin{foldable} % Our method TAKE and contributions
  We propose \textbf{Trajectory-Aware Knowledge Estimation} (\textbf{TAKE}), the first text \ac{DD} method to satisfy all three desiderata.
  TAKE corrects the \ac{HSB} by integrating influence scores along the full training trajectory into a per-sample knowledge score that captures clean and moderate samples while down-weighting noisy ones (\S\ref{sec:52-research03-theory}).
  Then, TAKE fine-tunes a \ac{LLM} to generate a pool of human-readable candidate instances.
  The knowledge scores and synthetic candidates jointly feed a discrete \ac{OT} objective that globally selects the distilled corpus.
  Although \ac{OT} has been widely applied in \ac{NLP} and machine learning, to our knowledge it has not been applied to text \ac{DD}, nor combined with a knowledge-reweighted source distribution in any \ac{DD} setting; TAKE fills both gaps.
  Concretely, we contribute:
  \begin{itemize}
    \item {
        \textbf{Theory:} We formalize the \ac{HSB} in the text \ac{DD} setting, prove that single-checkpoint influence scores are biased toward noisy samples, and derive a formal gap bound for knowledge-reweighted distribution matching.
      }
    \item {
        \textbf{Method:} TAKE is the first text \ac{DD} method to jointly address sample weighting and distributional coverage, via trajectory-integrated knowledge scores and a discrete \ac{OT} objective aligned with the \ac{DD} goal.
      }
    \item {
        \textbf{Empirical:} TAKE matches or exceeds prior text \ac{DD} methods across six language benchmarks at extreme compression (0.1\% or 20 instances/class), producing human-readable distilled corpora that generalize across diverse backbone families.
      }
  \end{itemize}
\end{foldable}

\begin{foldable} % Roadmap
  The chapter follows the order in which the three desiderata are settled.
  \S\ref{sec:52-research03-theory} casts text \ac{DD} as reweighted distribution matching problem and shows where the naive instantiation fails.
  \S\ref{sec:53-research03-method} then builds the TAKE pipeline: gradient-based trajectory influence and its knowledge score, synthetic candidate generation, and distributional prototype selection via discrete \ac{OT}.
  \S\ref{sec:54-research03-experiments} evaluates TAKE on classification and natural language inference benchmarks at extreme compression and confirms that each component is load-bearing.
  \S\ref{sec:55-research03-limitations} discusses limitations and ethical considerations, and \S\ref{sec:56-research03-conclusion} closes the chapter.
\end{foldable}

\section{Theoretical Framework} \label{sec:52-research03-theory}

\subsection{Reweighted Distribution Matching for Dataset Distillation} \label{sec:52-research03-theory-rdm}
\begin{foldable} % DD --> DM
    The \acl{DD} objective seeks $\tilde{\mathcal{D}}$ with $|\tilde{\mathcal{D}}| = K \ll N$ minimizing the loss on the original data:
    \begin{equation}
        \min_{\tilde{\mathcal{D}}}\; \mathbb{E}_{\theta_0 \sim p(\theta_0)} \left[ \mathbb{E}_{z \sim \mathcal{D}} \left[ \ell\!\left(z;\; F(\theta_0; \tilde{\mathcal{D}}, \eta, T)\right) \right] \right]
        \label{eq:dataset-distillation}
    \end{equation}
    where $F(\theta_0; \tilde{\mathcal{D}}, \eta, T)$ denotes the model obtained by running $T$ steps of gradient descent with learning rate $\eta$ on $\tilde{\mathcal{D}}$ from initialization $\theta_0$.
    This bi-level problem is intractable for discrete text.
    We relax it to distribution matching via the gradient-field equivalence argument~\cite{zhao2023dataset}: if $\tilde{\mathcal{D}}$ induces the same expected gradient field as $\mathcal{D}$ throughout training, the two runs converge to comparable parameters.
    Since $\nabla_\theta \mathbb{E}_{z \sim P}[\ell(z;\theta)] = \mathbb{E}_{z \sim P}[\nabla_\theta \ell(z;\theta)]$, matching distributions in a task-relevant feature space controls the gradient field --- making distribution matching a tractable surrogate for~\eqref{eq:dataset-distillation}.
    This yields the distribution matching~(DM) objective:
    \begin{equation}
        \min_{|\tilde{\mathcal{D}}| = K}\; d\!\left(P_{\tilde{\mathcal{D}}},\; P_{\mathcal{D}}\right)
        \label{eq:distribution-matching}
    \end{equation}
    for a divergence $d$ in a task-relevant feature space.
    Eq.~\eqref{eq:distribution-matching} is a surrogate, not an equivalence; the formal gap bound appears as Theorem~\ref{thm:gap} once $d$ is specified.
\end{foldable}

\begin{foldable} % Uniform DM --> RDM
    At extreme compression ($\rho = K/N \ll 1$), uniform matching allocates prototypes proportional to density --- but not all regions contribute equally to learning.
    Dense easy regions consume budget that would be better spent on moderate and boundary samples near decision boundaries, which carry the most information for model robustness.
    We instead reweight the source by per-sample importance: define $w_n \geq 0$, $\sum_n w_n = 1$, inducing $P^w = \sum_n w_n \delta_{z_n}$, yielding the \ac{RDM} objective:
    \begin{equation}
        \min_{|\tilde{\mathcal{D}}| = K}\; d\!\left(P_{\tilde{\mathcal{D}}},\; P^w\right)
        \label{eq:reweighted-distribution-matching}
    \end{equation}
    The weights $w_n$ must reflect downstream learning contribution, not geometry alone, motivating gradient-based estimation.
    Influence self-scores~\cite{koh2017understanding} are the natural candidate: they are task-aligned by construction and theoretically grounded.
    We now show that the simplest instantiation---$w_n \propto \mathcal{I}_\text{self}(z_n)$ at a single checkpoint, select top-$K$---fails on both counts: the weights are biased toward noisy samples (Proposition~\ref{prop:hsb}), and independent top-$K$ selection ignores distributional coverage (Remark~\ref{rem:coverage}).
    \begin{remark}[Double relaxation gap] \label{rem:double-gap}
        Moving from~\eqref{eq:dataset-distillation} to~\eqref{eq:distribution-matching} incurs a gradient-field approximation error (bounded formally by Theorem~\ref{thm:gap}); moving from~\eqref{eq:distribution-matching} to~\eqref{eq:reweighted-distribution-matching} replaces the uniform source $P_\mathcal{D}$ with the reweighted source $P^w$, incurring an additional cost:
        \begin{equation}
            \text{Gap}_2 \leq L \cdot W_1(P_\mathcal{D}, P^w) = L \cdot W_1\!\left(\tfrac{1}{N}\textstyle\sum_n \delta_{z_n},\; \textstyle\sum_n w_n \delta_{z_n}\right).
            \label{eq:gap2}
        \end{equation}
        Since both measures have the same support, $W_1(P_\mathcal{D}, P^w) \leq \text{diam}(\mathcal{Z}) \cdot \frac{1}{2}\sum_n |w_n - 1/N|$, which is directly computable from the $\kappa_n$ scores and small when $w_n \approx 1/N$ (uniform weights recover DM).
        The reweighting is beneficial precisely when this gap is outweighed by the reduction in downstream loss from correcting the \ac{HSB}, as formalized in Proposition~\ref{prop:traj} (\S\ref{sec:53-research03-method-gradient}).
    \end{remark}
\end{foldable}


\subsection{The Naive Instantiation Fails: Two Open Problems} \label{sec:52-research03-theory-failures}
\begin{foldable} % Setup
    Define the influence self-score at checkpoint $\hat{\theta}$ as
    \begin{equation}
        \mathcal{I}_\text{self}(z_n; \hat{\theta}) = \nabla_\theta \ell(z_n; \hat{\theta})^\top H_{\hat{\theta}}^{-1} \nabla_\theta \ell(z_n; \hat{\theta}),
        \label{eq:influence-self}
    \end{equation}
    where $H_{\hat{\theta}} = \frac{1}{N}\sum_n \nabla^2_\theta \ell(z_n; \hat{\theta})$ is the empirical Hessian, specializing~\cite{koh2017understanding} by setting the test point equal to the training point.
    High $\mathcal{I}_\text{self}$ identifies atypical examples: samples that are difficult to fit and exert outsized influence on the model's parameters, which in practice are the noisy ones.
    This suggests a naive pipeline: set $w_n \propto \mathcal{I}_\text{self}(z_n; \hat{\theta})$ and return $\arg\max_{|\mathcal{A}|=K} \sum_{i \in \mathcal{A}} w_i$; we identify two orthogonal failures of this strategy.
    The first failure is a bias in the weights themselves, due to the heterogeneous gradient magnitudes across clean and noisy samples that emerges over the training trajectory.
\end{foldable}

\begin{foldable} % Proposition 1: HSB
    \begin{proposition}[Hard-Sample Bias] \label{prop:hsb}
        Let $\mathcal{D}_\text{clean}, \mathcal{D}_\text{noisy}$ partition $\mathcal{D}$, and assume:
        \begin{enumerate}[label=\textbf{(A\arabic*)},leftmargin=*]
            \item $\|\nabla_\theta \ell(z; \hat{\theta})\| \leq \epsilon$ for $z \in \mathcal{D}_\text{clean}$, i.e., clean samples have near-zero gradient at convergence;
            \item $\|\nabla_\theta \ell(z; \hat{\theta})\| \geq \delta \gg \epsilon$ for $z \in \mathcal{D}_\text{noisy}$, i.e., noisy samples retain large residual gradients;
            \item $H_{\hat\theta}$ is symmetric positive definite with $\lambda_{\min}(H_{\hat\theta}) > 0$ (which holds whenever the model is not at a degenerate critical point).
        \end{enumerate}
        Let $\text{cond}(H) = \lambda_{\max}(H) / \lambda_{\min}(H)$ denotes the condition number of the Hessian $H$.
        Then:
        \begin{equation}
            \Delta_\textup{HSB} := \mathbb{E}_{z \sim \mathcal{D}_\textup{noisy}}[\mathcal{I}_\textup{self}(z;\hat{\theta})]
            - \mathbb{E}_{z \sim \mathcal{D}_\textup{clean}}[\mathcal{I}_\textup{self}(z;\hat{\theta})]
            \;\geq\; \frac{\delta^2 - \text{cond}(H_{\hat{\theta}}) \cdot \epsilon^2}{\lambda_{\max}(H_{\hat{\theta}})}
            \;>\; 0
            \label{eq:hsb}
        \end{equation}
    \end{proposition}
    \begin{proof}
        The argument bounds $\mathcal{I}_\text{self}(z)$ above for clean samples and below for noisy samples via the spectrum of $H_{\hat\theta}^{-1}$, then subtracts the resulting partition-wise expectations.
        Write $\nabla \equiv \nabla_\theta \ell(z;\hat\theta)$ so that $\mathcal{I}_\text{self}(z) = \nabla^\top H_{\hat\theta}^{-1} \nabla$.

        \textit{Step 1 (spectral bounds).}
        By (A3), $H_{\hat\theta}$ is symmetric positive definite, so $H_{\hat\theta}^{-1}$ is too, with eigenvalues in $[1/\lambda_{\max}(H_{\hat\theta}), 1/\lambda_{\min}(H_{\hat\theta})]$.
        Applying the Rayleigh quotient to $H_{\hat\theta}^{-1}$ gives:
        \begin{equation}
            \frac{\|\nabla\|^2}{\lambda_{\max}(H_{\hat\theta})}
            \;\leq\; \mathcal{I}_\text{self}(z)
            \;\leq\; \frac{\|\nabla\|^2}{\lambda_{\min}(H_{\hat\theta})}.
        \end{equation}
        Specializing to the two partitions using (A1) and (A2) and taking expectations (the bounds are uniform over each partition, so expectation preserves them):
        \begin{align}
            \mathbb{E}_{z \sim \mathcal{D}_\text{clean}}[\mathcal{I}_\text{self}(z)]
            &\leq \frac{\epsilon^2}{\lambda_{\min}(H_{\hat\theta})}, \\
            \mathbb{E}_{z \sim \mathcal{D}_\text{noisy}}[\mathcal{I}_\text{self}(z)]
            &\geq \frac{\delta^2}{\lambda_{\max}(H_{\hat\theta})}.
        \end{align}

        \textit{Step 2 (combining and positivity).}
        Subtracting the clean bound from the noisy bound, the difference of expectations satisfies
        \begin{equation}
            \mathbb{E}_{z \sim \mathcal{D}_\textup{noisy}}[\mathcal{I}_\textup{self}(z)] - \mathbb{E}_{z \sim \mathcal{D}_\textup{clean}}[\mathcal{I}_\textup{self}(z)]
            \;\geq\; \frac{\delta^2}{\lambda_{\max}(H_{\hat\theta})} - \frac{\epsilon^2}{\lambda_{\min}(H_{\hat\theta})}.
        \end{equation}
        Substituting $\lambda_{\min}(H_{\hat\theta}) = \lambda_{\max}(H_{\hat\theta})/\text{cond}(H_{\hat\theta})$ on the right-hand side and identifying the left-hand side with $\Delta_\textup{HSB}$ recovers \eqref{eq:hsb}.
        Positivity holds whenever $\delta/\epsilon > \sqrt{\text{cond}(H_{\hat\theta})}$, which is mild given $\delta \gg \epsilon$; in the interpolation regime ($\epsilon \to 0$) it is trivially satisfied.
    \end{proof}
    Because clean gradients vanish at convergence by definition of a local minimum, the bias is \textit{structural}: it cannot be corrected by post-hoc rescaling of single-checkpoint scores, but only by integrating gradient signals across the training trajectory.
    The binary partition simplifies exposition; in practice the bias is monotone in $\|\nabla\|$, as the spectral bounds above show directly.
    \begin{remark}[Gradient norm as special case] \label{rem:grad-norm}
        Setting $H = I$ in \eqref{eq:influence-self} recovers $\mathcal{I}_\text{self}(z) = \|\nabla\|^2$.
        Under this substitution, \eqref{eq:hsb} reduces to $\Delta_\text{HSB} \geq \delta^2 - \epsilon^2 > 0$, which holds without any spectral assumption on $H$.
        The gradient norm used in TAKE (\S\ref{sec:53-research03-method-gradient}) therefore inherits the same structural bias: it is not merely a tractable proxy but satisfies the \ac{HSB} lower bound directly.
    \end{remark}
\end{foldable}

\begin{foldable} % Remark 1: independent scoring != coverage
    Even with unbiased weights, a second failure remains: top-$K$ selection is blind to distributional coverage.
    \begin{remark}[Independent scoring $\neq$ distributional coverage] \label{rem:coverage}
        Any rank-and-select strategy decomposes as a sum of independent terms, invariant to pairwise distances between selected samples.
        In contrast, $d(P_{\tilde{\mathcal{D}}}, P^w)$ penalizes redundancy and rewards coverage.
        Concretely, if $\phi(z_i) \approx \phi(z_j)$ and both have high $w_i, w_j$, independent scoring selects both; a distributional objective assigns nearly the same cost to selecting just one, freeing a slot for an uncovered region.
    \end{remark}
\end{foldable}

\begin{foldable} % Two open problems
    The two failures compound: biased $w_n$ misdirects a population-level solver toward noisy regions; correct $w_n$ fed into an independent selector still produces redundancy.
    They respectively yield two open problems:
    \begin{enumerate}[label=\textbf{(\alph*)}]
        \item \label{open-problem-weight} \textbf{Biased weighting.} Estimate $w_n$ by integrating gradient signals across the learning trajectory $\{\theta_t\}_{t=0}^{T-1}$, reducing $\Delta_\textup{HSB}$---which we will address in \S\ref{sec:53-research03-method-gradient} (formalized as Proposition~\ref{prop:traj}).
        \item \label{open-problem-select} \textbf{Independent selection.} Solve \eqref{eq:reweighted-distribution-matching} directly for a metric $d$ that jointly optimises coverage of $P^w$---which we will address in \S\ref{sec:53-research03-method-ot}.
    \end{enumerate}
\end{foldable}
\section{Methodology} \label{sec:53-research03-method}

\subsection{Problem Setup and Framework Overview} \label{sec:53-research03-method-setup}
\begin{foldable} % Setup + overview merged
  TAKE instantiates the \ac{RDM} objective (Eq.~\ref{eq:reweighted-distribution-matching}) by solving the two open problems from \S\ref{sec:52-research03-theory-failures}.
  Let $\mathcal{D} = \{z_n\}_{n=1}^{N}$ denote the training corpus, $\mathcal{D}_\text{synth} = \{\tilde{z}_c\}_{c=1}^{M}$ the synthetic candidate pool, and $\tilde{\mathcal{D}} \subset \mathcal{D}_\text{synth}$ the distilled set of budget $K$.
  Indices $n \in [N]$ and $t \in [T]$ range over training samples and checkpoints respectively.
  TAKE proceeds in three stages:
  \begin{enumerate}
    \item \textbf{Score}---train a logistic probe on $\mathcal{D}$ across $T$ checkpoints, compute influence matrix $I \in \mathbb{R}^{N \times T}$, apply reciprocal reweighting and temporal weighting to obtain trajectory-aware knowledge scores $\kappa_n$.
    \item \textbf{Generate}---fine-tune a small language model on $\mathcal{D}$ and sample to produce the synthetic candidate pool $\mathcal{D}_\text{synth}$.
    \item \textbf{Distill}---solve discrete \ac{OT} with source weights $\kappa_n$ over $\mathcal{D}_\text{synth}$ to obtain the distilled corpus $\tilde{\mathcal{D}}$ of budget $K$.
  \end{enumerate}
  The three stages are decoupled and independently replaceable, with $\phi$ shared across Stages 1 and 3 as a natural consistency choice. The following sections detail each in turn.
\end{foldable}

\subsection{Gradient-Based Influence Along the Trajectory} \label{sec:53-research03-method-gradient}
\begin{foldable} % Motivation for trajectory-aware estimation
  The \acl{HSB} (\S\ref{sec:52-research03-theory-failures}) arises because single-checkpoint influence scores overemphasize noisy or difficult samples, whose gradients remain large at convergence.
  The fix is to integrate gradient information across the full training trajectory, producing a trajectory-aware knowledge score $\kappa_n$ per sample.
  We build $\kappa_n$ in two steps: constructing the influence matrix $I$ (\S\ref{sec:53-research03-method-gradient-influence}), then collapsing it into a per-sample scalar score via reciprocal reweighting and temporal weighting (\S\ref{sec:53-research03-method-gradient-knowledge}).
  The resulting $\kappa_n$ provides a task-aligned, bias-corrected weight for $P^w$.
\end{foldable}

\subsubsection{Influence Matrix} \label{sec:53-research03-method-gradient-influence}
\begin{foldable} % Definition and construction of I
  Our goal is to construct $I \in \mathbb{R}^{N \times T}$ whose entry $I_{(n,t)}$ reflects how much sample $z_n$ contributes to learning at checkpoint $\theta_t$.
  The influence self-score (\S\ref{sec:52-research03-theory-failures}, Eq.~\ref{eq:influence-self}) is the natural candidate, but requires Hessian inversion at a converged checkpoint and is undefined mid-trajectory.
  We address both limitations by adapting the per-sample gradient norm from TracIn~\cite{pruthi2020estimating} to the self-influence setting, querying training points instead of test points:
  \begin{equation}
    I_{(n,t)} = \|\nabla_\theta \ell(z_n; \theta_t)\|.
    \label{eq:influence-matrix}
  \end{equation}
  As the $H=I$ special case of the influence self-score (Remark~\ref{rem:grad-norm}), $I_{(n,t)}$ exhibits the same \ac{HSB} structure---the bias that $\kappa_n$ corrects via reciprocal inversion and trajectory integration (Proposition~\ref{prop:traj}).
  Scores are column-wise normalized $I_{(n,t)} \leftarrow I_{(n,t)} / \sum_{n'} I_{(n',t)}$ for comparability across checkpoints.
\end{foldable}

\begin{foldable} 
  In practice, we score samples with a linear head $W$ atop a pre-trained, frozen sentence encoder $\phi$, giving $\theta = (W, \phi)$.
  The signal quality comes not from the linear head but from $\phi(x_n)$: $\phi$ is pre-trained on large-scale corpora and $W$ is fine-tuned on the full task corpus, so the representations, and hence the gradient norms, are both semantically rich and task-aligned by construction.
  Following the standard practice in scalable influence estimation~\cite{guo2021fastif}, we approximate full-model influence via the last layer---here, the linear head $W$:
  \begin{equation}
    \nabla_\theta \ell(z_n; \theta_t) \approx \nabla_W \ell(z_n; \theta_t),
    \label{eq:last-layer}
  \end{equation}
  trading a small bias for tractability.
  By the Goodfellow outer-product identity~\cite{goodfellow2015efficient}, $\nabla_W \ell = g_n \phi(x_n)^\top$ where $g_n$ is the output gradient.
  The Goodfellow trick vectorizes the batch backward pass, avoiding per-sample loops.
  Adapting to a new task is trivial: retrain $W$ on the new corpus; $\phi$ and the influence pipeline are unchanged.
\end{foldable}

\subsubsection{Knowledge Score and Temporal Weighting} \label{sec:53-research03-method-gradient-knowledge}
\begin{foldable} % Inversion and temporal weighting to produce kappa
  To correct the \ac{HSB}, we require a score that is high for well-fitted samples and low for noisy ones---the opposite of what $I_{(n,t)}$ directly provides.
  Since a well-fitted sample has a small gradient norm, taking the reciprocal naturally inverts this ranking.
  We further weight each checkpoint by a decreasing temporal kernel $k : \{0,\ldots,T-1\} \to \mathbb{R}_{>0}$, giving:
  \begin{equation}
    \kappa_n = \sum_{t=0}^{T-1} \frac{k(t)}{I_{(n,t)}},
    \label{eq:kappa}
  \end{equation}
\end{foldable}

\begin{foldable}
  A decreasing kernel is the natural choice: early checkpoints capture easy samples fitting rapidly, mid checkpoints offer the strongest signal, and late checkpoints show narrowing as memorization sets in~\cite{arpit2017closer}.
  Thus, concentrating mass on the early-to-mid phase and discounting late memorization directly corrects the \ac{HSB} structure identified in Proposition~\ref{prop:hsb}.
  By default we use the exponential kernel $k(t) = e^{-\lambda t / T}$ ($\lambda > 0$), with $\lambda$ treated as a hyperparameter.
  Alternative kernels (linear, cosine) are drop-in replacements, and Proposition~\ref{prop:traj} guarantees bias reduction holds for any strictly decreasing $k$.
  The resulting $\kappa_n$ summarizes each sample's cumulative, bias-corrected learning contribution across the full trajectory, resolving open problem~\ref{open-problem-weight}.
  Normalized source weights $w_n = \kappa_n / \sum_{n'} \kappa_{n'}$ define the reweighted distribution $P^w$ and serve as input to the \ac{OT} selection step (\S\ref{sec:53-research03-method-ot}).
\end{foldable}

\begin{foldable}
  \begin{proposition}[Trajectory Integration Reduces HSB] \label{prop:traj}
    Let $T_m$ be the memorization onset, the first checkpoint at which noisy-sample gradient norms begin to decrease:
    \begin{equation*}
      T_m := \inf\{t : \mathbb{E}_{z \sim \mathcal{D}_\textup{noisy}}[\|\nabla_\theta \ell(z;\theta_t)\|] < \mathbb{E}_{z \sim \mathcal{D}_\textup{noisy}}[\|\nabla_\theta \ell(z;\theta_0)\|]\}.
    \end{equation*}
    Assume:
    \begin{enumerate}[label=\textbf{(A\arabic*)},leftmargin=*]
      \item \textit{Monotone clean decay.} For $z \in \mathcal{D}_\textup{clean}$, $I_{(n,t)}$ is non-increasing in $t$ and converges to zero (holds under $L$-smooth loss and a sufficiently small learning rate).
      \item \textit{Bounded noisy gradients up to memorization.} For $z \in \mathcal{D}_\textup{noisy}$, $I_{(n,t)} \geq \delta > 0$ for all $t \leq T_m$.
      \item \textit{Post-memorization baseline inflation.} At $t = T-1 > T_m$, noisy reciprocals exceed clean reciprocals, so $\Delta_{T-1} := \mathbb{E}_\textup{noisy}[1/I_{(n,T-1)}] - \mathbb{E}_\textup{clean}[1/I_{(n,T-1)}] > 0$.
      \item \textit{Positive decreasing kernel.} $k(t) > 0$ for all $t$, with $k(t^*) > k(T-1)$ for any $t^* < T-1$.
    \end{enumerate}
    Then $\kappa_n$ produces a strictly smaller noisy-over-clean gap than the single last-checkpoint baseline $\kappa^\dagger_n = 1/I_{(n,T-1)}$:
    \begin{equation}
      \mathbb{E}_{z \sim \mathcal{D}_\textup{noisy}}[\kappa_n] - \mathbb{E}_{z \sim \mathcal{D}_\textup{clean}}[\kappa_n]
      < \mathbb{E}_{z \sim \mathcal{D}_\textup{noisy}}[\kappa^\dagger_n] - \mathbb{E}_{z \sim \mathcal{D}_\textup{clean}}[\kappa^\dagger_n].
      \label{eq:prop2}
    \end{equation}
  \end{proposition}
  \begin{proof}
    The strategy is to track the per-checkpoint noisy-versus-clean gap $\Delta_t$ across the trajectory and show that the kernel-weighted sum on the left of \eqref{eq:prop2} is negative, while the last-checkpoint baseline on the right is positive.
    Define $\Delta_t := \mathbb{E}_\textup{noisy}[1/I_{(n,t)}] - \mathbb{E}_\textup{clean}[1/I_{(n,t)}]$, so the left-hand side of \eqref{eq:prop2} expands as:
    \begin{equation*}
      \mathbb{E}_\textup{noisy}[\kappa_n] - \mathbb{E}_\textup{clean}[\kappa_n] = \sum_{t=0}^{T-1} k(t) \cdot \Delta_t.
    \end{equation*}

    \textit{Step 1 (sign structure).}
    For $1 \leq t \leq T_m$: (A1) gives clean gradients decaying so $1/I_{(n,t)}$ grows for clean samples; (A2) bounds noisy gradients away from zero, so $1/I_{(n,t)} \leq 1/\delta$ for noisy samples.
    Once $t$ is late enough in the window $[1, T_m]$ that the clean reciprocal exceeds $1/\delta$ on average, $\Delta_t \leq 0$, strictly at some $t^*$ in that range; assumption (A1)'s convergence-to-zero clause guarantees such a $t^*$ exists.
    For $t > T_m$: $\Delta_t$ may be positive since noisy gradients decay after memorization.

    \textit{Step 2 (weighted sum).}
    Split the sum at $T_m$:
    \begin{equation*}
      S_1 = \sum_{t=1}^{T_m} k(t) \cdot \Delta_t \;<\; 0,
      \qquad
      S_2 = \sum_{t=T_m+1}^{T-1} k(t) \cdot \Delta_t \;\geq\; 0.
    \end{equation*}
    By (A4), $S_1 \leq -k(t^*) \cdot |\Delta_{t^*}|$ (single strictly negative term) and $S_2 \leq k(T_m+1) \cdot \sum_{t > T_m} \Delta_t$.
    Take the kernel $k(t) = e^{-\lambda^* t / T}$ with $\lambda^* = T/T_m$; the ratio $k(T_m+1)/k(t^*) = e^{-\lambda^*(T_m+1-t^*)/T}$ shrinks as $\lambda^*$ grows, so $S_2 / |S_1| \to 0$.
    For sufficiently large $\lambda^*$, $|S_1| > S_2$ and therefore $\sum_t k(t) \cdot \Delta_t = S_1 + S_2 < 0$.

    \textit{Step 3 (baseline gap).}
    By (A3), the single-checkpoint baseline satisfies $\mathbb{E}_\textup{noisy}[\kappa^\dagger_n] - \mathbb{E}_\textup{clean}[\kappa^\dagger_n] = \Delta_{T-1} > 0$.

    Combining Steps 2 and 3:
    \begin{equation*}
      \mathbb{E}_\textup{noisy}[\kappa_n] - \mathbb{E}_\textup{clean}[\kappa_n] \;<\; 0 \;<\; \Delta_{T-1} = \mathbb{E}_\textup{noisy}[\kappa^\dagger_n] - \mathbb{E}_\textup{clean}[\kappa^\dagger_n].
    \end{equation*}
  \end{proof}
  Operationally, this means trajectory-integrated scores favor clean and moderate samples over noisy ones, which is the ranking TAKE's selection step requires.
\end{foldable}

\subsection{Synthetic Candidate Pool Generation} \label{sec:53-research03-method-synthetic}

\begin{foldable} % Why a synthetic pool
  Rather than selecting directly from $\mathcal{D}$, TAKE draws $\tilde{\mathcal{D}}$ from a synthetic candidate pool $\mathcal{D}_\text{synth}$.
  Direct subsampling has two problems: at extreme compression, verbatim records raise privacy concerns, and the budget $K$ is too small to cover the full training support.
  We generate $\mathcal{D}_\text{synth}$ by fine-tuning a small language model on $\mathcal{D}$ and sampling from it.
  The fine-tuned LM generalizes beyond memorized training instances, and stochastic sampling produces a diverse pool that extends beyond the training support.
  $\mathcal{D}_\text{synth}$ serves solely as the selection pool and is never appended to downstream training data.
\end{foldable}

\begin{foldable} % Distribution mismatch is beneficial
  TAKE does \textit{not} assume $P_{\mathcal{D}_\text{synth}} = P_{\mathcal{D}}$; this distributional gap is \textit{beneficial}.
  First, coverage expansion fills undersampled regions of $\mathcal{D}$, improving generalization of the distilled corpus.
  Second, model fluency bias suppresses noisy samples that \ac{HSB} would otherwise over-select, yielding a cleaner candidate pool.
  Third, the distribution mismatch physically decouples $\tilde{\mathcal{D}}$ from original records, reducing verbatim exposure risk.
  \Acl{OT} handles the distributional gap gracefully: it routes mass toward $P^w$ regardless of pool distribution, so $P_{\tilde{\mathcal{D}}}$ approximates $P^w$ provided $\mathcal{D}_\text{synth}$ covers the support of $P_{\mathcal{D}}$ (Assumption~\ref{ass:coverage}):
  \begin{assumption}[Pool Coverage] \label{ass:coverage}
    For every training sample $z_n \in \mathcal{D}$, there exists a synthetic sample $\tilde{z}_c \in \mathcal{D}_\text{synth}$ within embedding distance $r$:
    $\min_c \|\phi(z_n) - \phi(\tilde{z}_c)\| \leq r$.
  \end{assumption}
  Under Assumption~\ref{ass:coverage}, $W_1(P^w, P_{\mathcal{D}_\text{synth}}) \leq r$, so the optimal transport cost is at most $r$ before any prototype selection.
  A practitioner can monitor coverage directly by reporting the mean nearest-neighbor distance $\bar{r} = \frac{1}{N}\sum_n \min_c \|\phi(z_n) - \phi(\tilde{z}_c)\|$.
\end{foldable}

\subsection{Distributional Prototype Selection} \label{sec:53-research03-method-ot}

\begin{foldable} % Why discrete OT
  Both $\mathcal{D}$ and $\mathcal{D}_\text{synth}$ are finite sample sets, so the matching problem is naturally over empirical measures.
  Discrete \ac{OT} operates directly on $\mu = \sum_n w_n \delta_{z_n}$ and $\nu = \frac{1}{|\mathcal{D}_\text{synth}|}\sum_c \delta_{\tilde{z}_c}$, requiring no density estimation and yielding an explicit transport plan with direct prototype-level assignments.
\end{foldable}

\begin{foldable} % Formulation and entropic regularization
  Let $C_{nc} = \|\phi(z_n) - \phi(\tilde{z}_c)\|^2$ be the cost matrix in a task-relevant embedding space.
  Bare Kantorovich \ac{OT}, $\min_{\gamma \in \Pi(\mu,\nu)} \langle C, \gamma \rangle$, produces a sparse, non-differentiable plan and scales as $O(N^3)$.
  We add entropic regularization to obtain a convex, smooth objective solvable via Sinkhorn--Knopp iterations at $O(N \cdot |\mathcal{D}_\text{synth}|)$ per step:
  \begin{equation}
    \gamma^* = \arg\min_{\gamma \in \Pi(\mu,\nu)} \langle C, \gamma \rangle + \varepsilon \cdot \text{KL}(\gamma \| \mu \otimes \nu).
    \label{eq:ot}
  \end{equation}
  The soft coupling distributes mass across nearby candidates rather than hard one-to-one assignment, reducing sensitivity to embedding errors and encouraging diversity in $\tilde{\mathcal{D}}$.
  As an extension, the independence reference $\mu \otimes \nu$ can be replaced with a task-specific prior coupling $\alpha(z, \tilde{z})$ (\eg, a topic-conditional measure) to regularize the transport plan toward domain knowledge.
  The $K$ candidates with highest received mass $\sum_n \gamma^*_{nc}$ are selected as $\tilde{\mathcal{D}}$, resolving open problem~\ref{open-problem-select}.
  The following theorem closes the theoretical chain from \ac{RDM} (Eq.~\ref{eq:reweighted-distribution-matching}) back to \ac{DD} (Eq.~\ref{eq:dataset-distillation}), bounding the downstream loss gap in terms of the \ac{OT} cost.

  \begin{theorem}[Distribution Matching Gap Bound] \label{thm:gap}
    Let $f^*_{\mathcal{S}}$ denote the model trained on $\mathcal{S}$, and let $\mathcal{L}(\mathcal{S}) = \mathbb{E}_{z \sim P^w}[\ell(z; f^*_{\mathcal{S}})]$.
    Assume:
    \begin{enumerate}[label=\textbf{(A\arabic*)},leftmargin=*]
      \item $\ell(\cdot; f)$ is $L$-Lipschitz in the embedding metric for all $f$;
      \item the learning algorithm satisfies uniform stability: $\sup_z |\ell(z; f^*_{\tilde{\mathcal{D}}}) - \ell(z; f^*_\mathcal{D})| \leq \beta_K$ with $\beta_K = O(1/K)$~\cite{bousquet2002stability};
      \item the ground metric for $W_1$ is the embedding distance $c(z, z') = \|\phi(z) - \phi(z')\|$.
    \end{enumerate}
    Then:
    \begin{equation}
      \bigl|\mathcal{L}(\tilde{\mathcal{D}}) - \mathcal{L}(\mathcal{D})\bigr|
      \;\leq\; 2L \cdot W_1(P^w, P_{\tilde{\mathcal{D}}}) + \beta_K,
      \label{eq:gap-bound}
    \end{equation}
    where $W_1$ is the Wasserstein-1 distance under the embedding metric.
    As $K \to \infty$, $\beta_K \to 0$, and the bound is dominated by the $W_1$ term controlled by the \ac{OT} objective.
  \end{theorem}
  \begin{proof}
    The proof inserts the surrogate distribution $P_{\tilde{\mathcal{D}}}$ as a triangulation pivot between $P^w$ and the model-difference term, then bounds the two distribution-shift pieces by Wasserstein duality and the model-difference piece by uniform stability.

    \textit{Step 1 (Kantorovich--Rubinstein bound for a fixed model).}
    By Kantorovich--Rubinstein duality~\cite{villani2008optimal} under (A3), for any fixed $f$:
    \begin{equation*}
      \bigl|\mathbb{E}_{z \sim P^w}[\ell(z; f)] - \mathbb{E}_{z \sim P_{\tilde{\mathcal{D}}}}[\ell(z; f)]\bigr|
      \;\leq\; L \cdot W_1(P^w, P_{\tilde{\mathcal{D}}}).
    \end{equation*}

    \textit{Step 2 (triangle inequality decomposition).}
    $\mathcal{L}(\tilde{\mathcal{D}})$ uses $f^*_{\tilde{\mathcal{D}}}$ while $\mathcal{L}(\mathcal{D})$ uses $f^*_\mathcal{D}$; both are evaluated under $P^w$ but with different models, so Step 1 does not apply directly.
    Inserting two intermediates and applying the triangle inequality twice:
    \begin{align*}
      |\mathcal{L}(\tilde{\mathcal{D}}) - \mathcal{L}(\mathcal{D})|
      &\leq \underbrace{\bigl|\mathbb{E}_{P^w}[\ell(z;f^*_{\tilde{\mathcal{D}}})] - \mathbb{E}_{P_{\tilde{\mathcal{D}}}}[\ell(z;f^*_{\tilde{\mathcal{D}}})]\bigr|}_{\text{(I) distribution shift, fixed }f^*_{\tilde{\mathcal{D}}}} \\
      &\quad + \underbrace{\bigl|\mathbb{E}_{P_{\tilde{\mathcal{D}}}}[\ell(z;f^*_{\tilde{\mathcal{D}}})] - \mathbb{E}_{P_{\tilde{\mathcal{D}}}}[\ell(z;f^*_\mathcal{D})]\bigr|}_{\text{(II) model difference, fixed }P_{\tilde{\mathcal{D}}}} \\
      &\quad + \underbrace{\bigl|\mathbb{E}_{P_{\tilde{\mathcal{D}}}}[\ell(z;f^*_\mathcal{D})] - \mathbb{E}_{P^w}[\ell(z;f^*_\mathcal{D})]\bigr|}_{\text{(III) distribution shift, fixed }f^*_\mathcal{D}}.
    \end{align*}
    Terms (I) and (III) each satisfy Step 1 (with $f = f^*_{\tilde{\mathcal{D}}}$ and $f = f^*_\mathcal{D}$ respectively), so each is bounded by $L \cdot W_1(P^w, P_{\tilde{\mathcal{D}}})$.
    Term (II) is bounded pointwise by (A2): $\text{(II)} \leq \beta_K = O(1/K)$.

    Combining the three: $|\mathcal{L}(\tilde{\mathcal{D}}) - \mathcal{L}(\mathcal{D})| \leq 2L \cdot W_1(P^w, P_{\tilde{\mathcal{D}}}) + \beta_K$.
  \end{proof}
  The bound is stated with respect to $P^w$ rather than the uniform $P_\mathcal{D}$ --- this is intentional, as $P^w$ is the designed optimization target of the \ac{RDM} objective; since $P^w$ reweights $P_\mathcal{D}$ over the same support, minimizing $W_1(P^w, P_{\tilde{\mathcal{D}}})$ controls coverage of $P_\mathcal{D}$ up to the reweighting gap bounded in Remark~\ref{rem:double-gap}.
\end{foldable}

\begin{foldable} % Method wrap-up
  Together, Propositions~\ref{prop:hsb}--\ref{prop:traj} and Theorem~\ref{thm:gap} establish the full theoretical basis for \ac{TAKE}: \ac{HSB} is structural and unavoidable at a single checkpoint (Proposition~\ref{prop:hsb}); trajectory integration strictly reduces it (Proposition~\ref{prop:traj}); and minimizing the \ac{OT} cost directly controls the downstream loss gap under $P^w$ (Theorem~\ref{thm:gap}; the gap to $P_\mathcal{D}$ is bounded by Remark~\ref{rem:double-gap}).
  Algorithm~\ref{alg:take} consolidates the full \ac{TAKE} pipeline.
\end{foldable}

\begin{center}
  \begin{minipage}{0.99\linewidth}
    \begin{algorithm}[H]
      \SetKwInOut{KwIn}{Input}
      \SetKwInOut{KwOut}{Output}
      \SetKwComment{Comment}{}{}

      \caption{Trajectory-Aware Knowledge Estimation (\textbf{TAKE}).}
      \label{alg:take}

      \KwIn{corpus $\mathcal{D} = \{z_n\}_{n=1}^{N}$; budget $K$; checkpoints $T$; kernel $k(\cdot)$; $\varepsilon$; pool size $M$}
      \KwOut{distilled corpus $\tilde{\mathcal{D}}$, $|\tilde{\mathcal{D}}| = K$}

      \vspace{-0.5\baselineskip}
      \Comment{\noindent\hrulefill}

      \vspace{0.5\baselineskip}
      \Comment{$\triangleright$ \textbf{\textit{Stage 1: Score}}}
      Train logistic probe on $\mathcal{D}$, saving checkpoints $\theta_0, \ldots, \theta_{T-1}$ \\
      \For{$t = 0, \ldots, T-1$}{
        Compute $I_{(n,t)} \leftarrow \|\nabla_W \ell(z_n; \theta_t)\|$ for all $n$ \hfill\textit{(Goodfellow trick)} \\
        Normalize $I_{(\cdot,t)} \leftarrow I_{(\cdot,t)} / \sum_{n'} I_{(n',t)}$ \hfill\textit{(column-wise)}
      }
      Compute $\kappa_n \leftarrow \textstyle\sum_{t=0}^{T-1} k(t) / I_{(n,t)}$ for all $n$ \hfill via Eq.~\ref{eq:kappa} \\
      Normalize $w_n \leftarrow \kappa_n / \sum_{n'} \kappa_{n'}$ for all $n$ \hfill\textit{(defines $P^w$)}

      \vspace{0.5\baselineskip}
      \Comment{$\triangleright$ \textbf{\textit{Stage 2: Generate}}}
      Fine-tune small language model $\text{LM}$ on $\mathcal{D}$ \\
      Sample synthetic pool $\mathcal{D}_\text{synth} = \{\tilde{z}_c\}_{c=1}^{M}$ from $\text{LM}$

      \vspace{0.5\baselineskip}
      \Comment{$\triangleright$ \textbf{\textit{Stage 3: Distill}}}
      Embed: $\phi(z_n)$ for $n \in [N]$; $\phi(\tilde{z}_c)$ for $c \in [M]$ \\
      Compute cost matrix $C_{nc} = \|\phi(z_n) - \phi(\tilde{z}_c)\|^2$ \\
      Solve $\gamma^* \leftarrow \text{Sinkhorn}(C, w, \mathbf{1}/M, \varepsilon)$ \hfill via Eq.~\ref{eq:ot} \\
      Distill $\tilde{\mathcal{D}} \leftarrow \text{top}\text{-}K\!\left(\{\tilde{z}_c\},\; \textstyle\sum_n \gamma^*_{nc}\right)$ \hfill\textit{(by received mass)}

      \vspace{0.5\baselineskip}
      \Return{$\tilde{\mathcal{D}}$}
    \end{algorithm}
  \end{minipage}
\end{center}

\input{./content/50-research03/54-research03-experiments}
\section{Limitations and Ethical Considerations} \label{sec:55-research03-limitations}

\begin{foldable} % Limitations
  This method has three limitations.
  (\textbf{1})~\textbf{Gradient norm proxy.}
  $I_{(n,t)}$ uses plain gradient norms rather than Fisher-preconditioned scores; the Goodfellow trick~\cite{goodfellow2015efficient} keeps per-sample cost at $O(C + d_\phi)$ per checkpoint.
  K-FAC~\cite{martens2015optimizing} or EKFAC~\cite{george2018fast} preconditioning are extensions to consider when calibrated magnitude scores are needed.
  (\textbf{2})~\textbf{Encoder dependence.}
  The \ac{OT} cost matrix inherits the geometry of $\phi$; a mismatched encoder distorts transport costs.
  The $L$-Lipschitz assumption on $\ell$ (Theorem~1) holds approximately when $\phi$ is task-aligned and embeddings are normalized, but remains an approximation for cross-entropy over discrete text.
  (\textbf{3})~\textbf{Synthetic pool coverage.}
  Pool diversity is bounded by the generator's domain coverage.
  Fine-tuning on $\mathcal{D}_\text{train}$ helps address this.
\end{foldable}

\begin{foldable} % Ethical considerations
  Two ethical considerations apply to TAKE.
  (\textbf{1})~\textbf{Bias propagation.}
  Distillation concentrates the statistical properties of $\mathcal{D}_\text{train}$, including demographic biases.
  Corpus auditing before distillation is recommended.
  (\textbf{2})~\textbf{Privacy risk.}
  Language models fine-tuned on sensitive data may memorize training examples, a risk not fully eliminated by synthetic data generation.
  Distance-to-closest-record (DCR) screening of $\mathcal{D}_\text{synth}$ can verify that synthetic samples are not verbatim copies; for regulated domains, differential privacy at the fine-tuning stage provides a formal guarantee against memorization.
\end{foldable}

\section{Conclusion} \label{sec:56-research03-conclusion}

\begin{foldable} % Conclusion
  We propose TAKE, the first text \ac{DD} method to simultaneously satisfy all three desiderata: \textit{knowledge-based weighting} via trajectory-integrated influence scores that correct the \acl{HSB} and direct the budget toward informative samples; a \textit{dataset-level distillation objective} via discrete Optimal Transport that globally aligns the selected corpus with the source distribution; and \textit{human-readable output} via selection from \ac{LLM}-generated candidates, producing a distilled corpus that is inspectable, auditable, and transferable across architectures.
  These design choices address gaps that prior text \ac{DD} methods leave open: TAKE matches or exceeds baselines across six language benchmarks at extreme compression, with end-to-end theoretical formulation and guarantees.
  These results suggest that text dataset distillation can help address systemic cost pressures: reducing corpus size cuts cumulative training expenditure, lowers carbon footprint, and respects data-governance constraints without sacrificing transferability.
\end{foldable}

